\documentclass[11pt,a4paper]{article}

%% PACHETE MATEMATICE
\usepackage{amsmath,amssymb,amsfonts}
\usepackage{amsthm}
\usepackage{mathpazo}
\usepackage{eulervm}
\usepackage{enumerate}
\usepackage{color}
\usepackage{graphicx}
\usepackage[all]{xy}
\usepackage{xcolor}
\usepackage{framed}
\usepackage[unicode]{hyperref}
\setcounter{MaxMatrixCols}{20}
\usepackage{multicol}
% \usepackage[romanian]{babel}


%% ASPECT
\linespread{1.05}
\usepackage[T1]{fontenc}
\usepackage[utf8x]{inputenc}
\usepackage[margin=2cm]{geometry}
% \usepackage[color]{showkeys}
\usepackage{anyfontsize}
\usepackage{titlesec}
\titleformat*{\section}{\large\bfseries}

\usepackage{fancyhdr}
\pagestyle{fancy}
\rhead{\small \color{gray} Notițe de Adrian Manea}
\lhead{\small \color{gray} IntProb-FIA, 2017---2018, L. Lemnete \& M. Olteanu}


\usepackage[autostyle = false, style = german]{csquotes}
\usepackage[romanian]{babel}
\newcommand\qq{\enquote}

%% COMENZILE MELE
\renewcommand*{\proofname}{Dem.:}
\newcommand\bb{\mathbb}
\newcommand\dr{\mathrm}
\newcommand\kal{\mathcal}
\newcommand\fk{\mathfrak}
\newcommand\op{\oplus}
\newcommand\ot{\otimes}
\newcommand\ds{\displaystyle}
\newcommand\id{\indent}
\newcommand\nid{\noindent}
\newcommand\seq{\subseteq}
\newcommand\sneq{\subsetneq}
\newcommand\speq{\supseteq}
\newcommand\spneq{\supsetneq}
\newcommand\rar{\rightarrow}
\newcommand\Rar{\Rightarrow}
\newcommand\xrar{\xrightarrow}
\newcommand\lar{\leftarrow}
\newcommand\Lar{\Leftarrow}
\newcommand\xlar{\xleftarrow}
\newcommand\lrar{\leftrightarrow}
\newcommand\Lrar{\Leftrightarrow}
\newcommand\ideal{\trianglelefteq}
\newcommand\ai{\text{ a.\^{i}. }}
\newcommand\sk{(\textit{Schi\c{t}\u{a}}) }
\newcommand\rhup{\rightharpoonup}
\newcommand\rhdn{\rightharpoondown}
\newcommand\lhup{\leftharpoonup}
\newcommand\lhdn{\leftharpoondown}
\newcommand\vid{\emptyset}
\newcommand\wh{\widehat}
\newcommand\ol{\overline}
\newcommand\NN{\mathbb{N}}
\newcommand\RR{\mathbb{R}}
\newcommand\ZZ{\mathbb{Z}}
\newcommand\QQ{\mathbb{Q}}
\newcommand\CC{\mathbb{C}}

%% STILURILE MELE
\newtheoremstyle{dotless}{}{}{\itshape}{}{\bfseries}{:}{ }{}

\theoremstyle{dotless}
\newtheorem{theorem}{Teorem\u{a}}[section]
\newtheorem{proposition}{Propozi\c{t}ie}[section]
\newtheorem{lemma}{Lem\u{a}}[section]
\newtheorem{corollary}{Corolar}[section]
\newtheorem{exercise}{Exerci\c{t}iu}

\newtheoremstyle{dotlessdr}{}{}{}{}{\bfseries}{:}{ }{}
\theoremstyle{dotlessdr}
\newtheorem{example}{Exemplu}[section]
\newtheorem{definition}{Defini\c{t}ie}[section]
\newtheorem{remark}{Observa\c{t}ie}[section]




\begin{document}

\begin{center}
  {\large\textbf{Seminar 7 \\ Integrale duble și triple}}
\end{center}

\vspace{1cm}

1. Calculați integralele duble:
\begin{enumerate}[(a)]
\item $ \ds\iint_D ydxdy $, unde $ D $ este mărginit de parabola $ y^2 = 2x $,
  cercul $ x^2 + y^2 - 2x = 0 $ și dreapta $ x = 2 $;
\item $ \ds\iint_D e^{x^2 + y^2} dxdy $, unde $ D = \{(x, y) \in \RR^2 \mid x^2 + y^2 \leq 1 \} $;
\item $ \ds\iint_D ydxdy $, unde $ D = \{(x, y) \in \RR^2 \mid (x - 2)^2 + y^2 \leq 1 \} $.
\end{enumerate}

\vspace{1cm}

2. Fie $ D \seq \RR^2 $ și $ f : D \to [0, \infty) $ o funcție continuă. Definim
mulțimea:
\[
  \Omega = \{ (x, y, z) \in \RR^3 \mid (x, y) \in D, 0 \leq z \leq f(x, y) \}.
\]
Să se calculeze volumul mulțimii $ \Omega $, folosind formula:
\[
  \kal{V}(\Omega) = \iint_D f(x, y) dxdy,
\]
în cazurile:
\begin{enumerate}[(a)]
\item $ D = \{(x, y) \in \RR^2 \mid x^2 + y^2 \leq 2y \}, f(x, y) = x^2 + y^2 $;
\item $ D = \{(x, y) \in \RR^2 \mid x^2 + y^2 \leq x, y > 0 \}, f(x, y) = xy $;
\item $ D = \{(x, y) \in \RR^2 \mid x^2 + y^2 \leq 2x + 2y - 1 \}, f(x, y) = y $.
\end{enumerate}

\vspace{1cm}

\emph{Indicații:} (a) Mulțimea este $ x^2 + (y - 1)^2 \leq 1 $, care este discul
centrat în $ (0, 1) $ și de rază 1. Putem folosi coordonate polare pentru a face
calculul, remarcînd că, deoarece discul se află deasupra axei $ OX $, avem
$ \theta \in [0, \pi] $.

În plus, din inegalitatea $ x^2 + y^2 \leq 2y $, obținem $ \rho \leq 2 \sin \theta $.

Atunci integrala este:
\[
  \int_0^\pi d\theta \int_0^{2 \sin \theta} \rho^3 d\rho = \frac{3}{2}\pi.
\]

\vspace{1cm}

3. Să se calculeze integrala triplă:
\[
  I = \iiint_V \sqrt{x^2 + y^2} dxdydz,
\]
unde $ V $ este delimitat de suprafața conică $ z = \sqrt{x^2 + y^2} $ și planul
$ z = a > 0 $.

\vspace{1cm}

4. Să se calculeze integrala triplă:
\[
  I = \iiint_V xydxdydz,
\]
unde $ V $ este mărginit de cilindrul $ x^2 + y^2 = R^2 $ și planele
$ z = 0, z = 1, y = x, y = x\sqrt{3} $.

\vspace{1cm}

5. Să se calculeze volumul corpului delimitat de suprafețele:
\[
  \begin{cases}
    z &= x^2 + y^2 \\
    z &= 2x^2 + 2y^2 \\
    y &= x \\
    y &= x^2
  \end{cases}.
\]

\vspace{1cm}

6. Să se calculeze volumul cuprins între paraboloizii:
\[
  \begin{cases}
    x^2 + y^2 &= 2 + z \\
    x^2 + y^2 &= 10 - z
  \end{cases}.
\]

\vspace{1cm}

7. Calculați volumele mulțimilor $ \Omega $, mărginite de suprafețele de ecuații:
\begin{enumerate}[(a)]
\item $ 2x^2 + y^2 + z^2 = 1, \quad 2x^2 + y^2 - z^2 = 0, z \geq 0 $;
\item $ z = x^2 + y^2 - 1, \quad z = 2 - x^2 - y^2 $;
\item $ z = 4 - x^2 - y^2, \quad 2z = 5 + x^2 + y^2 $;
\item $ x^2 + y^2 = 1, \quad z^2 = x^2 + y^2, \quad z \geq 0 $.
\end{enumerate}

\emph{Indicații:} (a) Intersectăm cele două suprafețe, egalînd primele două ecuații și
obținem:
\[
  4x^2 + 2y^2 = 1, \quad z = \frac{1}{\sqrt{2}}.
\]
Atunci, proiecția acestui domeniu pe planul $ XOY $ este:
\[
  D = \{(x, y) \mid 4x^2 + 2y^2 \leq 1 \}.
\]
Iar coordonata $ z $ o putem lua între cele două suprafețe, adică:
\[
  \sqrt{2x^2 + y^2} \leq z \leq \sqrt{1 - 2x^2 - y^2}.
\]
Rezultă că putem calcula volumul mulțimii astfel:
\[
  V(\Omega) = \iiint_\Omega dxdydz = \iint_D dxdy \int_{\sqrt{2x^2 + y^2}}^{\sqrt{1 - 2x^2 - y^2}} dz,
\]
iar apoi pe domeniul $ D $ putem folosi coordonate polare generalizate (pentru parametrizarea
elipsei).

\vspace{1cm}

(b) Similar, intersectăm cele două suprafețe rezolvînd sistemul de ecuații
și găsim:
\[
  x^2 + y^2 = \frac{3}{2}, \quad z = \frac{1}{2}.
\]
Proiecția pe planul $ XOY $ este:
\[
  D = \big\{ (x, y) \mid x^2 + y^2 \leq \frac{3}{2} \big\}.
\]
Atunci:
\[
  V(\Omega) = \iiint_\Omega dxdydz = \iint_D dxdy \int_{x^2 + y^2 - 1}^{2 - x^2 - y^2} dz.
\]




\end{document}